\documentclass[10pt,a4paper]{amsart}
\usepackage[margin=19mm]{geometry}
\usepackage{amsmath,amssymb,mathtools}
\usepackage[scr=rsfs,cal=euler]{mathalfa}
\usepackage{xfrac}
\usepackage[unicode,hidelinks,bookmarks=false]{hyperref}
\newcommand{\ie}{i.e.\ }
\newcommand{\cf}{cf.\ }
\newcommand{\resp}{respectively\ }
\newcommand{\Subsection}{\subsection}
\providecommand{\nobreakdash}{\nobreak\hbox{-}}
\setlength{\emergencystretch}{2em}
\multlinegap0pt
\usepackage{bm}
\usepackage{bbm}
\usepackage{dsfont}
\usepackage{xspace}
\usepackage{cancel}
\usepackage{mathtools}
\usepackage{amsthm}
\makeatletter
\@ifundefined{lem}{\newtheorem{lem}{Lemma}}{}
\@ifundefined{thm}{\newtheorem{thm}[lem]{Theorem}}{}
\makeatother
\newcommand{\N}{\mathbb{N}}
\newcommand{\R}{\mathbb{R}}
\newcommand{\norm}[1]{\left\Vert#1\right\Vert}
\renewcommand{\geq}{\geqslant}
\renewcommand{\leq}{\leqslant}
\title[$C^1$ invariant, stable and inertial manifolds for non-auton]{$C^1$ invariant, stable and inertial manifolds for non-autonomous dynamical systems}
\author{Rados{\l}aw Czaja, Piotr Kalita, Alexandre N. Oliveira-Sousa}
\date{Source: arXiv:2508.00165v3 (CC BY 4.0)}
\begin{document}
\maketitle

\noindent\emph{Source and licence.} $C^1$ invariant, stable and inertial manifolds for non-autonomous dynamical systems. Version arXiv:2508.00165v3 (CC BY 4.0), distributed under the Creative Commons Attribution 4.0 International licence, as recorded in the arXiv licence metadata for this exact version and shown on the abstract page https://arxiv.org/abs/2508.00165v3. Full rights notice: licenses/arxiv-2508.00165-CC-BY-4.0.txt.\par\smallskip

\noindent\textit{Editorial scope.} Counted unit: the Thm whose source label is \texttt{e:CONTINFSIGMA} in the article (source0.tex lines 1755--1762, source label \texttt{e:CONTINFSIGMA}), delivered together with the article's own complete original proof of it (source0.tex lines 1764--1855, one \texttt{proof} environment of 92 lines). No mathematics is written, repaired or re-derived: statement and proof are the article's own text line for line. Required context retained uncounted: e:QI-QONX (lines 212--214); lem:EQUIVNORMS (lines 234--241); e:L1DF (lines 1536--1538); e:L2DF (lines 1539--1541); e:CONTINUITYPSIT (lines 1568--1574); e:DEFPSI (lines 1641--1643); e:FRECHET3 (lines 1739--1741). Every definition, display and label of those lines is retained unchanged. Omitted from this package and not redistributed: 1--211, 215--233, 242--1535, 1542--1567, 1575--1640, 1644--1738, 1742--1754, 1763--1763, 1856--1936. Nothing inside a retained excerpt is omitted or altered: every retained body is delivered exactly as the article's lines are written, so no omission receipt of any kind accompanies this package. Citations: the retained excerpts carry no citation command at all, so no bibliography entry of the article is transcribed and no reference is redistributed. Reference adaptation: none was needed and none was made; every reference inside the delivered unit resolves inside it, so no unresolved reference or citation is produced. Printed slips kept literal: no printed word, symbol, hypothesis or number of the counted statement or of its proof was altered. Layout and class: the article's own class is \texttt{amsart} with packages including amsmath, amsfonts, amsthm, amssymb, dsfont, amsaddr, comment, enumitem, hyperref, graphicx, color, srcltx; the wrapper is the lane's pinned \texttt{amsart} editorial wrapper at 10pt on A4, which restates the theorem-like environments the retained text needs and adds the article's own macro definitions that the retained text needs (\texttt{\textbackslash N}, \texttt{\textbackslash R}, \texttt{\textbackslash geq}, \texttt{\textbackslash leq}, \texttt{\textbackslash norm}), with extra wrapper packages (bm, bbm, dsfont, xspace, cancel, mathtools). The delivered numbering is the unit's own. Assets: no retained span contains a figure environment, a graphics-inclusion command, a PDF-page inclusion, a table or a TikZ picture; no image file is copied into this package, the package declares no asset dependency and no image file is redistributed with it. Licence: version arXiv:2508.00165v3 of this article is distributed under the Creative Commons Attribution 4.0 International licence, as recorded in the arXiv licence metadata for this exact version and shown on the abstract page https://arxiv.org/abs/2508.00165v3. A full rights notice is delivered at licenses/arxiv-2508.00165-CC-BY-4.0.txt. Characters: every retained excerpt is pure ASCII, so the delivered engine, XeLaTeX with its default Latin Modern text font, reports no missing character.
\begin{equation}\label{e:QI-QONX}
|Q(\tau)x|_{N(\tau)}\leq M\norm{x}\ \text{and}\  |(I-Q(\tau))x|_{S(\tau)}\leq M\norm{x}\ \text{for}\  x\in X.
\end{equation}

\begin{lem}\label{lem:EQUIVNORMS}
Given an admissible norm on $\R^2$, for each $\tau\in\R$ the function
\begin{equation}\label{e:MOVINGNORM}
\norm{x}_\tau=\Gamma\left(|Q(\tau)x|_{N(\tau)},|(I-Q(\tau))x|_{S(\tau)}\right)\ \text{for}\ x\in X
\end{equation}
is an equivalent norm in $X$. We have 
$$\frac{1}{c_{\Gamma}}\norm{x}\leq\norm{x}_{\tau}\leq M\Gamma(1,1)\norm{x}\text{  for }x\in X\ \text{and}\ \tau\in\R.$$
\end{lem}

\begin{equation}\label{e:L1DF}
	|Q(t)Df(t,u)w|_{N(t)}\leq L_1\norm{w}_{t},
\end{equation}

\begin{equation}\label{e:L2DF}
	|(I-Q(t))Df(t,u)w|_{S(t)}\leq L_2\norm{w}_{t}.
\end{equation}

\begin{align}\label{e:PSIONNT}
&|Q(t)[\psi(\eta+h)(t)-\psi(\eta)(t)]|_{N(t)}\leq|Q(\tau)h|_{N(\tau)}e^{(\rho+L_1\Gamma(1,\kappa_\Sigma))(\tau-t)},\\
&\label{e:PSIONST}
|(I-Q(t))[\psi(\eta+h)(t)-\psi(\eta)(t)]|_{S(t)}\leq\kappa_\Sigma|Q(\tau)h|_{N(\tau)}e^{(\rho+L_1\Gamma(1,\kappa_\Sigma))(\tau-t)},\\
&\label{e:CONTINUITYPSIT}
\norm{\psi(\eta+h)(t)-\psi(\eta)(t)}_{t}\leq\Gamma\left(1,\kappa_\Sigma\right)|Q(\tau)h|_{N(\tau)}e^{(\rho+L_1\Gamma(1,\kappa_\Sigma))(\tau-t)}.
\end{align}

\begin{equation}\label{e:DEFPSI}
\Psi(\eta)=Z_{\eta}\ \ \text{for}\ \ \eta\in X.
\end{equation}

\begin{equation}\label{e:FRECHET3}
\norm{Df(t,u)-Df(t,u_0)}_{\mathcal{L}(X,X)}\to 0\text{ as }u\to u_0\text{ in }X\text{ for every }t\in\R,\ u_0\in X. 
\end{equation}

\begin{thm}
The function $X\ni \eta\mapsto \Psi(\eta)\in F_\sigma\subset F_\mu$ defined in \eqref{e:DEFPSI} is continuous in $F_\mu$, i.e.,
$$\norm{\Psi(\eta+h)-\Psi(\eta)}_{F_\mu}\to 0\text{ as }h\to0\text{ for each }\eta\in X.$$ 
In particular, we have
\begin{equation}\label{e:CONTINFSIGMA}
\norm{\Psi(\eta+h)(t)-\Psi(\eta)(t)}_{\mathcal{L}(X,X)}\to0\text{ as }h\to0\ \text{for each}\ t\leq\tau\ \text{and}\ \eta\in X.
\end{equation}
\end{thm}

\begin{proof}
We first observe that $\Psi$ defined in \eqref{e:DEFPSI} satisfies
\begin{align*}
	& \Psi(\eta+h)(t)-\Psi(\eta)(t)=-\int_{t}^{\tau}L(t,s)Q(s)Df(s,\psi(\eta+h)(s))\Psi(\eta+h)(s)ds\\
	&\ \ \ \  +\int_{t}^{\tau}L(t,s)Q(s)Df(s,\psi(\eta)(s))\Psi(\eta)(s)ds\\
	&\ \ \ \  +\int_{-\infty}^{t}L(t,s)(I-Q(s))Df(s,\psi(\eta+h)(s))\Psi(\eta+h)(s)ds\\
& \ \ \ \  -\int_{-\infty}^{t}L(t,s)(I-Q(s))Df(s,\psi(\eta)(s))\Psi(\eta)(s)ds\ \ \text{for}\ \ t\leq\tau,\ \eta, h\in X.
\end{align*}
Thus, we have
\begin{align*}
	& \Psi(\eta+h)(t)-\Psi(\eta)(t)=-\int_{t}^{\tau}L(t,s)Q(s)[Df(s,\psi(\eta+h)(s))-Df(s,\psi(\eta)(s))]\Psi(\eta)(s)ds\\
	& \ \ -\int_{t}^{\tau}L(t,s)Q(s)Df(s,\psi(\eta+h)(s))[\Psi(\eta+h)(s)-\Psi(\eta)(s)]ds\\
& \ \ +\int_{-\infty}^{t}L(t,s)(I-Q(s))[Df(s,\psi(\eta+h)(s))-Df(s,\psi(\eta)(s))]\Psi(\eta)(s)ds\\
& \ \ +\int_{-\infty}^{t}L(t,s)(I-Q(s))Df(s,\psi(\eta+h)(s))[\Psi(\eta+h)(s)-\Psi(\eta)(s)]ds\ \ \text{for}\ \ t\leq\tau,\ \eta, h\in X.
\end{align*}
Hence, by \eqref{e:L1DF}, for $\norm{w}\leq 1$ we have 
\begin{align*}
&	|Q(t)[\Psi(\eta+h)(t)-\Psi(\eta)(t)]w|_{N(t)}\leq L_1\int_{t}^{\tau}e^{\rho(s-t)}\norm{[\Psi(\eta+h)(s)-\Psi(\eta)(s)]w}_{s}ds \\
& \ \  +\int_{t}^{\tau}e^{\rho(s-t)}|Q(s)[Df(s,\psi(\eta+h)(s))-Df(s,\psi(\eta)(s))]\Psi(\eta)(s)w|_{N(s)}ds.
\end{align*}
Similarly, by \eqref{e:L2DF}, for $\norm{w}\leq 1$ we have 
\begin{align*}
	&|(I-Q(t))[\Psi(\eta+h)(t)-\Psi(\eta)(t)]w|_{S(t)}\leq L_2\int_{-\infty}^{t}e^{\gamma(s-t)}\norm{[\Psi(\eta+h)(s)-\Psi(\eta)(s)]w}_{s}ds\\
	& \ \ +\int_{-\infty}^{t}e^{\gamma(s-t)}|(I-Q(s))[Df(s,\psi(\eta+h)(s))-Df(s,\psi(\eta)(s))]\Psi(\eta)(s)w|_{S(s)}ds.\end{align*}
Thus, we get
\begin{align*}
	&e^{\mu(t-\tau)}|Q(t)[\Psi(\eta+h)(t)-\Psi(\eta)(t)]w|_{N(t)}\leq L_1\norm{\Psi(\eta+h)-\Psi(\eta)}_{F_\mu}\int_{t}^{\tau}e^{(\mu-\rho)(t-s)}ds\\
	& \ \ +e^{\mu(t-\tau)}\int_{t}^{\tau}e^{\rho(s-t)}|Q(s)[Df(s,\psi(\eta+h)(s))-Df(s,\psi(\eta)(s))]\Psi(\eta)(s)w|_{N(s)}ds,\\
& e^{\mu(t-\tau)}|(I-Q(t))[\Psi(\eta+h)(t)-\Psi(\eta)(t)]w|_{S(t)}\leq L_2\norm{\Psi(\eta+h)-\Psi(\eta)}_{F_\mu}\int_{-\infty}^{t}e^{(\gamma-\mu)(s-t)}ds\\& \ \ +e^{\mu(t-\tau)}\int_{-\infty}^{t}e^{\gamma(s-t)}|(I-Q(s))[Df(s,\psi(\eta+h)(s))-Df(s,\psi(\eta)(s))]\Psi(\eta)(s)w|_{S(s)}ds.\end{align*}
Therefore, we obtain
$$\norm{\Psi(\eta+h)-\Psi(\eta)}_{F_\mu}\leq\Gamma\left(\frac{L_1}{\mu-\rho},\frac{L_2}{\gamma-\mu}\right)\norm{\Psi(\eta+h)-\Psi(\eta)}_{F_\mu}+\Gamma(A_1(\eta,h),B_1(\eta,h)),$$
where
\begin{align*}
	&A_1(\eta,h) =\sup_{t\leq\tau}\sup_{\norm{w}\leq 1}\Big(e^{(\mu-\rho)(t-\tau)}\cdot\\
	& \qquad \qquad  \qquad  \cdot \int_{t}^{\tau}e^{\rho(s-\tau)}|Q(s)[Df(s,\psi(\eta+h)(s))-Df(s,\psi(\eta)(s))]\Psi(\eta)(s)w|_{N(s)}ds\Big),\\
	& B_1(\eta,h) =\sup_{t\leq\tau}\sup_{\norm{w}\leq 1}\Big(e^{(\gamma-\mu)(\tau-t)}\cdot\\
	& \qquad \qquad  \qquad  \cdot\int_{-\infty}^{t}e^{\gamma(s-\tau)}|(I-Q(s))[Df(s,\psi(\eta+h)(s))-Df(s,\psi(\eta)(s))]\Psi(\eta)(s)w|_{S(s)}ds\Big).\end{align*}
Since $\Gamma\left(\frac{L_1}{\mu-\rho},\frac{L_2}{\gamma-\mu}\right)<1$, we will prove the claim if we show that
%$$(i)\quad A_1(\eta,h)\to 0\text{ as }h\to 0,$$
%$$(ii)\quad B_1(\eta,h)\to 0\text{ as }h\to 0.$$
$${A_1(\eta,h)\to 0\ \text{ and }\ B_1(\eta,h)\to 0\ \text{ as }\ h\to 0.}$$
To this end, first note that $A_1(\eta,h)$ and $B_1(\eta,h)$ are well defined. Indeed, by \eqref{e:L1DF}, \eqref{e:L2DF} we have
for $t\leq\tau$ and $\norm{w}\leq 1$
\begin{align*}
&	e^{(\mu-\rho)(t-\tau)}\int_{t}^{\tau}e^{\rho(s-\tau)}|Q(s)[Df(s,\psi(\eta+h)(s))-Df(s,\psi(\eta)(s))]\Psi(\eta)(s)w|_{N(s)}ds\\
	& \ \ \ \ \leq2L_1\norm{\Psi(\eta)}_{F_\sigma}e^{(\mu-\rho)(t-\tau)}\int_{t}^{\tau}e^{(\sigma-\rho)(\tau-s)}ds\leq\frac{2L_1}{\sigma-\rho}\norm{\Psi(\eta)}_{F_\sigma}e^{(\mu-\sigma)(t-\tau)},\\
& e^{(\gamma-\mu)(\tau-t)}\int_{-\infty}^{t}e^{\gamma(s-\tau)}|(I-Q(s))[Df(s,\psi(\eta+h)(s))-Df(s,\psi(\eta)(s))]\Psi(\eta)(s)w|_{S(s)}ds\\
& \ \ \ \ \leq 2L_2\norm{\Psi(\eta)}_{F_\sigma}e^{(\gamma-\mu)(\tau-t)}\int_{-\infty}^{t}e^{(\gamma-\sigma)(s-\tau)}ds=\frac{2L_2}{\gamma-\sigma}\norm{\Psi(\eta)}_{F_\sigma}e^{(\mu-\sigma)(t-\tau)}.
\end{align*}
Thus we have
$$A_1(\eta,h)\leq\frac{2L_1}{\sigma-\rho}\norm{\Psi(\eta)}_{F_\sigma}\ \text{ and }\ B_1(\eta,h)\leq\frac{2L_2}{\gamma-\sigma}\norm{\Psi(\eta)}_{F_\sigma}.$$
Suppose, contrary to the claim, that $A_1(\eta,h)\not\to 0$ as $h\to 0$. Thus there exist $\epsilon_0>0$ and $h_n\to 0$ such that $A_1(\eta,h_n)>\epsilon_0$ for all $n\in\N$. Therefore, there exist $t_n\leq\tau$ and $\norm{w_n}\leq 1$ such that 
$$\epsilon_0<e^{(\mu-\rho)(t_n-\tau)}\int_{t_n}^{\tau}e^{\rho(s-\tau)}|Q(s)[Df(s,\psi(\eta+h_n)(s))-Df(s,\psi(\eta)(s))]\Psi(\eta)(s)w_n|_{N(s)}ds.$$
Since the right-hand side is bounded from above by $\frac{2L_1}{\sigma-\rho}\norm{\Psi(\eta)}_{F_\sigma}e^{(\mu-\sigma)(t_n-\tau)}$ and $\mu>\sigma$,
the sequence $t_n$ cannot have a subsequence tending to $-\infty$. Hence there exists $\tau_0$ such that $t_n\geq\tau_0$ for all $n\in\N$.
We obtain
\begin{equation}\label{e:Q100}
\epsilon_0<\int_{\tau_0}^{\tau}e^{\mu(s-\tau)}|Q(s)[Df(s,\psi(\eta+h_n)(s))-Df(s,\psi(\eta)(s))]\Psi(\eta)(s)w_n|_{N(s)}ds.
\end{equation}
For each  $s\in[\tau_0,\tau]$ we have 
$$e^{\mu(s-\tau)}|Q(s)[Df(s,\psi(\eta+h_n)(s))-Df(s,\psi(\eta)(s))]\Psi(\eta)(s)w_n|_{N(s)}\leq 2L_1\norm{\Psi(\eta)}_{F_\mu},$$
which is integrable on $[\tau_0,\tau]$.

Now, recalling \eqref{e:QI-QONX} 
%that
%$$|Q(s)x|_{N(s)}\leq M\norm{x}\text{ and }|(I-Q(s))x|_{S(s)}\leq M\norm{x}\  \text{for}\  \ s\in\R,\ x\in X.$$
and deducing from Lemma~\ref{lem:EQUIVNORMS} that for $G\in\mathcal{L}(X,X)$
$$\norm{Gx}\leq c_\Gamma\norm{G}_{\mathcal{L}(X,X)}\norm{x}_{s},\ s\in\R,\ x\in X,$$
we further have
\begin{align*}
&e^{\mu(s-\tau)}|Q(s)[Df(s,\psi(\eta+h_n)(s))-Df(s,\psi(\eta)(s))]\Psi(\eta)(s)w_n|_{N(s)}\\
	&\qquad \qquad  \leq Mc_\Gamma\norm{Df(s,\psi(\eta+h_n)(s))-Df(s,\psi(\eta)(s))}_{\mathcal{L}(X,X)}\norm{\Psi(\eta)}_{F_\mu}.
\end{align*}
By \eqref{e:CONTINUITYPSIT} and \eqref{e:FRECHET3} this shows that the integrand on the right-hand side of \eqref{e:Q100} tends to zero as $n\to\infty$ pointwise on $[\tau_0,\tau]$. 
Thus the Lebesgue dominated convergence theorem shows that the right-hand side of \eqref{e:Q100} converges to zero as $n\to\infty$. This is a contradiction with the choice of $\epsilon_0$.

In a similar way, suppose now, contrary to the claim, that $B_1(\eta,h)\not\to 0$ as $h\to 0$. Thus there exist $\epsilon_0>0$ and $h_n\to 0$ such that $B_1(\eta,h_n)>\epsilon_0$ for all $n\in\N$. Therefore, there exist $t_n\leq\tau$ and $\norm{w_n}\leq 1$ such that 
$$\epsilon_0<e^{(\gamma-\mu)(\tau-t_n)}\int_{-\infty}^{t_n}e^{\gamma(s-\tau)}|(I-Q(s))[Df(s,\psi(\eta+h_n)(s))-Df(s,\psi(\eta)(s))]\Psi(\eta)(s)w_n|_{S(s)}ds.$$
Since $(\gamma-\mu)(\tau-t_n)+\gamma(s-\tau)=\mu(s-\tau)+(\gamma-\mu)(s-t_n)$, we get
\begin{equation}\label{e:Q101}
\epsilon_0<\int_{-\infty}^{\tau}e^{\mu(s-\tau)}|(I-Q(s))[Df(s,\psi(\eta+h_n)(s))-Df(s,\psi(\eta)(s))]\Psi(\eta)(s)w_n|_{S(s)}ds.
\end{equation}
Now, for every $s\in(-\infty,\tau]$ we have 
$$e^{\mu(s-\tau)}|(I-Q(s))[Df(s,\psi(\eta+h_n)(s))-Df(s,\psi(\eta)(s))]\Psi(\eta)(s)w_n|_{S(s)}\leq 2L_2\norm{\Psi(\eta)}_{F_\sigma}e^{(\mu-\sigma)(s-\tau)},$$
which is an integrable function on $(-\infty,\tau]$, since $\mu>\sigma$. Moreover, we get 
\begin{align*}
	&e^{\mu(s-\tau)}|(I-Q(s))[Df(s,\psi(\eta+h_n)(s))-Df(s,\psi(\eta)(s))]\Psi(\eta)(s)w_n|_{S(s)}\\
	& \ \ \leq Mc_\Gamma\norm{Df(s,\psi(\eta+h_n)(s))-Df(s,\psi(\eta)(s))}_{\mathcal{L}(X,X)}\norm{\Psi(\eta)}_{F_\mu}.
\end{align*}
By \eqref{e:CONTINUITYPSIT} and \eqref{e:FRECHET3} this shows that the integrand on the right-hand side of \eqref{e:Q101} tends to zero as $n\to\infty$ pointwise on $(-\infty,\tau]$. 
Thus the Lebesgue dominated convergence theorem shows that the right-hand side of \eqref{e:Q101} tends to zero as $n\to\infty$. This is a contradiction with the choice of $\epsilon_0$.
\end{proof}

\end{document}
